\begin{textsample}{POS Dim 4 – human – Score 4.00 – mg/mg\_ef\_linguagens\_ingles\_12\_p1.txt}  \label{ex:f4_pos_015}
Ganham importância os estágios do processo cíclico de produção \textbf{textual} : geração de ideias ( brainstorming ) e reconhecimento pela \textbf{leitura} da organização \textbf{textual} do mesmo gênero que será produzido, planejamento, múltiplos rascunhos, reescritas, edição final, sempre com o suporte de opções diferentes de feedback fornecido pelos colegas, amigos e professor.
Subjacente está a noção de aperfeiçoamento do \textbf{texto} ao longo do processo de discussões, reflexões, rascunhos sucessivos e reescritas até a versão final, que inclui também a incorporação de recursos gráficos ( diagramação da página impressa, saliências de ênfase : negrito, itálico, tipologias diferenciadas ; ilustrações ; diagramas ; figuras ; legendas ; tabelas, etc. ), incluindo ainda preocupações com relação aos portadores de \textbf{textos} ( jornais, revistas, brochuras, Internet, livros, etc. ), que também exercem influência no processo de produção \textbf{textual}.

% matched lemmas: leitura, texto, textual
\end{textsample}
